\section{Introduction}
\label{sec:introduction}

Mass spectrometry imaging (MSI) maps molecular measurements directly to tissue
coordinates, enabling biochemical structure to be studied without preselecting
one target molecule~\citep{caprioli1997maldi,schwamborn2010molecular}. Each
pixel nevertheless contains a spectrum with hundreds to tens of thousands of
mass-to-charge ($m/z$) bins. Many bins are redundant, noisy or unrelated to the
biological structure of interest, so selecting a compact set of informative
peaks is an important preprocessing step.

Conventional peak picking uses rules such as intensity and signal-to-noise
thresholds. These rules are useful but can be sensitive to user choices and may
discard weak yet spatially coherent ions. Peak learning instead trains a model
on spectra and derives a ranking from the learned representation. The msiPL
approach uses a one-dimensional variational autoencoder (VAE), but processes
each pixel independently~\citep{abdelmoula2021peak}. S3PL explicitly consumes a
spatial patch~\citep{weigand2026spatial}, although its three-dimensional
convolutional architecture and released implementation differ substantially
from msiPL. This leaves two questions entangled: whether local context improves
the learned representation, and whether a better explanation of a nonlinear
representation improves the resulting peak list.

This report develops Spatial-msiPL, a controlled VAE study in which the central
spectrum is compared with the same model supplied with a summary of its Moore
neighbourhood. Rather than assuming that first-layer weights faithfully explain
the model, a label-free Gaussian mixture model (GMM) is fitted in latent space
and its posterior is propagated to input bins with Integrated
Gradients~\citep{sundararajan2017axiomatic}. Expert masks are used only after
ranking to quantify spatial peak quality.

The study answers three research questions:
\begin{enumerate}
  \item Under matched training, how does eight-neighbour context affect
  reconstruction, latent tissue separation and matched-count peak quality
  relative to a centre-only VAE?
  \item Does nonlinear GMM-targeted Integrated Gradients produce better peak
  rankings than legacy msiPL and simple first-layer L2 importance?
  \item What peak-quality, explanation-faithfulness and computational trade-offs
  arise from learned attention, and are they stable across training seeds?
\end{enumerate}

The aim is to determine separately when neighbourhood context and nonlinear
attribution improve lightweight MSI peak learning. The principal contributions
are: (i) a controlled centre-only/contextual comparison; (ii) a label-free
GMM-targeted attribution pipeline; (iii) matched-count validation on eight GBM
and eight CAC sections; and (iv) completeness, deletion-faithfulness, compute
and training-seed analyses that bound the claims. The remainder of the report
reviews related peak-learning and attribution methods, defines the controlled
pipeline, presents the experimental results, and discusses what they do and do
not establish.
